\section{Limitations and Discussion}
\label{sec:limitations}

The theory has sharp but restrictive scope. Target-interface agreement is an identification assumption, not an observable finite-sample certificate. The transport result separately requires a one-step coupling and Lipschitz continuation laws, which can fail at discontinuous safety-filter active-set changes. Small Wasserstein error does not imply stable high quantiles without local distributional regularity. In the current deterministic simulator, physical tail-risk language would therefore be misleading.

The empirical claim is also conditional. If exact Resource-to-Go has no headroom over simple heuristics, the problem may not warrant a learned manager in this domain. If executed-WM plus Deep Ensemble explains failures, SIRP should be deleted. If only the primary UAV simulator exhibits the mechanism, the conclusion must remain domain-specific. Finally, the learned module does not guarantee charger reachability or collision avoidance; it supports a stopping decision while HOCBF retains collision-safety authority.

These limitations suggest a broader interpretation. The useful abstraction is not “battery prediction,” but a stopping problem in which a resource forecast is evaluated through the interface actually executed and through the decision boundary it can cross. The same analysis may apply to robots returning for charge, tools, data upload, or thermal recovery, but such transfer is a hypothesis until replicated.
